\section{结论}
\label{sec:conclusion}

本文把告警分诊流程从小型本地 LLM 中移出，交给 \hesp{}：一个把调查拆成“下一步探什么”和“证据何时已足够”的控制器。带有计数似然表的假设账本按每单位成本的信息增益为只读探针排序，证据规则和控制器侧停止管控结论，预写日志使每一步都可审计。用五个开源权重模型、7{,}272 个回合完成的四项预注册研究表明：小模型在流程上失败；“选探针”和“决定停止”是两种独立的失败，不同的小模型表现各异；把流程从模型手中拿走，可以把 Qwen2.5-7B 的已验证完成率从 0.125 提升到 1.000，把 Llama-3.1-8B 从 0 提升到 0.917。完全不含 LLM 的控制器同样达到 0.917。在每个起因都留下独特离散观测的环境里，起作用的是流程；由此得到的结论关乎控制权而不是能力：小型本地模型不应掌控调查流程。至于模型在最需要它的地方——阅读模糊的原始遥测、区分看起来像日常工作的攻击和看起来像攻击的日常工作——能否带来价值，需要本文以及我们考察过的公开安全数据集都不具备的环境和数据。
